\chapter{Estados visuales}
\label{anexo:estados-visuales}

Resumen para consulta durante el torneo. Para interpretacion completa, consulte \cref{cap:estados-mensajes}.

\begin{longtable}{p{4.0cm}p{6.6cm}p{4.3cm}}
\toprule
Estado visual & Significado & Accion recomendada\\
\midrule
\endhead
Pendiente & Falta completar o guardar un dato. & Complete el flujo relacionado.\\
Seleccionado & El operador marco una opcion visible. & Guarde antes de avanzar.\\
Modificado & Hay cambios respecto al estado inicial de la pantalla. & Use guardado individual o general.\\
Guardado & La operacion fue aceptada por el servidor. & Recargue o revise persistencia.\\
Activo & Participante sigue en competencia. & Mantenga en la fase correspondiente.\\
Clasificado & Participante avanza desde grupo o batalla. & Revise fase siguiente.\\
Ganador & Participante gano una batalla o final. & Compare con resultado oficial.\\
Eliminado o perdedor & Participante no avanza desde ese resultado. & No corregir sin autorizacion.\\
En repechaje & Participante puede competir en fase de retorno. & Revise candidatos antes de crear repechaje.\\
Retirado & Participante fue marcado fuera de competencia. & Confirmar decision antes de operar.\\
Incompleto & Faltan campos, resultados o confirmaciones. & No avanzar hasta cerrar pendientes.\\
Bloqueado & La interfaz impide continuar. & Lea mensaje y vuelva al paso anterior.\\
Finalizado & Batalla, fase o competencia cerrada. & Continue o cierre segun corresponda.\\
Borrador & Comunicacion o competencia aun no finalizada. & Revisar antes de enviar o avanzar.\\
Previsualizado & Comunicacion simulada o plantilla revisada. & Confirmar antes de envio real.\\
En envio & Lote en proceso de envio. & Espere historial; no duplique lote.\\
Enviado & Correo real enviado. & Auditar destinatarios e historial.\\
Prueba enviada & Correo enviado a receptor de prueba. & Validar contenido antes de envio real.\\
Simulado & No hubo envio fisico. & Usar para validacion inicial.\\
Omitido & Registro no procesado por condicion o error de tipo. & Revisar causa antes de repetir.\\
Fallido & La operacion no se completo. & Corregir causa y escalar si persiste.\\
Cancelado & Proceso detenido o sin registros utiles. & Verificar si debe repetirse.\\
\bottomrule
\end{longtable}
